笛卡爾座標直接表示原子在空間中的位置，建模介面通常以埃（\AA）為單位。分數座標則表示沿三條晶格向量的比例。若分數座標為 $(u,v,w)$，原子的位置就是 $u\mathbf a+v\mathbf b+w\mathbf c$。這三個數沒有長度單位；它們回答的是“沿每條晶格向量走多少”，而不是“沿全域性座標軸走多少埃”。

為讀懂晶格輸入表格，可以把三條向量按行排成一個矩陣：
\[
 A=\begin{bmatrix}
 a_x&a_y&a_z\\
 b_x&b_y&b_z\\
 c_x&c_y&c_z
 \end{bmatrix},\qquad
 \mathbf f=(u,v,w),\qquad \mathbf r=(x,y,z).
\]
矩陣的第一行是 $\mathbf a$，第二、三行依次是 $\mathbf b$、$\mathbf c$；每行中的三列才是全域性 $x,y,z$ 分量。本書將座標寫成行向量，因此兩種座標的轉換為
\begin{equation}
 \mathbf r=\mathbf f A,\qquad \mathbf f=\mathbf r A^{-1}.
 \label{eq:row-lattice-coordinates}
\end{equation}
這裡 $A^{-1}$ 是整個晶格矩陣的逆矩陣，不能把表格中的每個數分別取倒數。晶格三行所定義的晶胞應有非零體積，才能由笛卡爾位置求回唯一的分數座標。

以一組便於手算的晶格說明：令 $\mathbf a=(4,0,0)$、$\mathbf b=(1,3,0)$、$\mathbf c=(0,0,6)$，單位均為\AA。分數座標 $(0.5,0.5,0)$ 對應
\[
 \mathbf r=0.5(4,0,0)+0.5(1,3,0)
          =(2.5,1.5,0)\,\text{\AA}.
\]
沿 $\mathbf b$ 走一半，同時包含全域性 $x$ 方向的 $0.5\,\text{\AA}$ 和 $y$ 方向的 $1.5\,\text{\AA}$。所以斜晶胞中的第二個分數分量，不能直接當作全域性 $y$ 座標。圖\ref{fig:round10-fractional-cartesian}將這一次相加畫成兩段相接的箭頭。

\begin{figure}[!htb]
\centering
\begin{tikzpicture}[x=1mm,y=1mm,font=\sffamily\small]
 \fill[paleblue] (8,6)--(56,6)--(68,42)--(20,42)--cycle;
 \draw[teal,line width=.7pt] (8,6)--(56,6)--(68,42)--(20,42)--cycle;
 \draw[->,>=stealth,muted] (8,6)--(84,6) node[right] {$x$};
 \draw[->,>=stealth,muted] (8,6)--(8,49) node[above] {$y$};
 \draw[->,>=stealth,brand,line width=1pt] (8,6)--(32,6);
 \draw[->,>=stealth,brand,line width=1pt] (32,6)--(38,24);
 \draw[->,>=stealth,teal,line width=.8pt] (8,6)--(38,24);
 \draw[dashed,muted] (38,24)--(38,6);
 \fill[brand] (38,24) circle (1.2);
 \node[anchor=north,text=brand] at (22,5) {$u\mathbf a$};
 \node[anchor=west,text=brand] at (38,14) {$v\mathbf b$};
 \node[anchor=south east,text=teal] at (25,16) {$\mathbf r$};
 \node[anchor=north east] at (7,5) {$O$};
 \node[anchor=west,text=brand] at (42,27) {$\mathbf f=(u,v,0)$};
 \node[anchor=west,text=muted,font=\sffamily\footnotesize] at (78,35) {晶胞在平面内的截面};
 \node[anchor=west,text=muted,font=\sffamily\footnotesize] at (78,28) {箭头相加得到位置};
 \node[anchor=north,text=teal] at (54,4) {$\mathbf a$};
 \node[anchor=east,text=teal] at (15,31) {$\mathbf b$};
\end{tikzpicture}
\caption{分数坐标与笛卡尔位置的关系（二维示意）。$u\mathbf a$ 与 $v\mathbf b$ 相加得到 $\mathbf r$；斜晶胞的 $\mathbf b$ 方向同时含有全局 $x$ 和 $y$ 分量。}
\label{fig:round10-fractional-cartesian}
\end{figure}


在邊長為 $5\,\text{\AA}$ 的立方晶胞中，$(0.5,0,0)$ 則對應沿 $\mathbf a$ 前進 $2.5\,\text{\AA}$。兩種例子的區別來自晶格，而不是座標轉換規則不同。編輯時若保持 $\mathbf f$ 不變而改變 $A$，原子通常會隨晶胞移動；若希望原子在空間中的位置不動，則需保留 $\mathbf r$，並用新晶格重新求 $\mathbf f$。第\ref{ch:modeling-tools}章的\ref{sec:lattice-transform-reference}節介紹這兩個選項。

\begin{table}[!htb]
\centering\small
\begin{tabularx}{\linewidth}{llX}
\toprule
符號 & 單位 & 在引數中的含義\\
\midrule
$A$ & \AA & 晶格矩陣，三行依次為 $\mathbf a,\mathbf b,\mathbf c$。\\
$\mathbf f$ & 無量綱 & 沿三條晶格向量的比例，即分數座標。\\
$\mathbf r,\mathbf p$ & \AA & 笛卡爾位置，以及旋轉錨點或鏡面上的點。\\
$\mathbf n,\hat{\mathbf n}$ & 方向 & 非零方向向量及其單位向量；只看方向時不以\AA 填長度。\\
$\theta$ & 度 & 對話方塊中的轉角，正角按右手定則。\\
$S$ & 整數 & 定義新晶格的超胞變換矩陣。\\
$N,V$ & 個、\AA$^3$ & 模型獨立位點數與晶胞體積。\\
$\varepsilon_a,\varepsilon_b,\varepsilon_c$ & 無量綱 & 沿三條晶格向量的相對伸縮量。\\
\bottomrule
\end{tabularx}
\caption{閱讀建模公式時常用的符號。座標和錨點使用長度，方向向量只指定朝向；兩者在對話方塊中有不同作用。}
\label{tab:round10-modeling-symbols}
\end{table}
